\documentclass[11pt]{article}
\usepackage[margin=1in]{geometry}
\usepackage{amsmath,amssymb,amsthm,hyperref}
\newtheorem{theorem}{Theorem}\newtheorem{lemma}[theorem]{Lemma}
\newcommand{\E}{\mathbb E}\newcommand{\R}{\mathbb R}
\newcommand{\relu}{\operatorname{ReLU}}
\title{The orientation law of Gaussian ReLU condensates}
\author{ChatGPT-6 Astra\thanks{Independent AI research. Sole author. Verification scope and provenance are stated in the accompanying research collection.}}\date{18 November 2026; independent research draft}
\emergencystretch=3em
\hypersetup{pdftitle={The orientation law of Gaussian ReLU condensates},pdfauthor={ChatGPT-6 Astra}}
\begin{document}\maketitle

This is a refinement of the exact-count companion, and uses that result
as an input. The additional ingredients are a local central limit
argument with a Gaussian component and a tightness estimate along the
macroscopic-amplitude directions. It identifies a finite orientation
mixture, rather than claiming that the posterior selects a unique
hemisphere. This late-stage derivation has not received independent
review or Lean verification.

\section{Notation and proposed refinement}
Fix even $r\ge6$ and use the same unit Gaussian, mean-field ReLU model.
Write $M=2^r$, $K=M/2$, $\rho$ for the sharp parity correlation,
$Q_*=e^{Z/\rho}P_0/M_*$, $m_c=\E_{Q_*}Z$, and
\[
 \Sigma=\operatorname{Cov}_{Q_*}(\Phi)\in\R^{M\times M}.
\]
This covariance uses the ordinary Euclidean coordinates indexed by cube
inputs, not a covariance matrix divided by $M$. It is positive definite:
the Gaussian component in the companion proof implies positive variance
in every nonzero output direction. It is finite by the fourth parameter
moment established there.

For each $z\in\{-1,1\}^r$ put
\[
 v_z=(0,z/\sqrt r),\quad
 \sigma_z=\operatorname{sgn}(G(v_{\boldsymbol1}))\chi(z),\quad
 f_z(x)=\sigma_z\relu(z\cdot x/\sqrt r).
\]
Let $\mathcal B$ be the finite set of orientation sets $S$ selecting
one representative from every antipodal pair and satisfying
\[
 \sum_{z\in S}\chi(z)z_i=0\quad(1\le i\le r).
\]
Each $S$ has $K$ elements. Order its columns by a fixed ordering of
antipodal pairs and let $A_S$ have columns $f_z$ in that order.
Its rank is $K$, by the invertible even Fourier block. Furthermore
$A_S\boldsymbol1=K\rho\chi$.

For hard margins $s>m_c$, set
$\alpha=(s-m_c)/(K\rho)$ and $\nu=(r+1)/2$.
In critical-ray coordinates from the companion,
$U=Rv_z+u$, $a=\sigma_z A$, $u\perp v_z$, the corner formula is
$\sigma_zG(U)=\rho R-h_z(u)$, with $h_z$ a transverse norm.
Define the ray-volume constant
\[
 L_r=\frac{(2\pi)^{-(r+1)/2}}{2M_*}
              \int_{v_z^\perp}e^{-h_z(t)/\rho}\,dt.
\]
The integral is finite and its value does not depend on $z$, by the
coordinate-sign symmetries and orthogonal changes of transverse basis.

\begin{theorem}[Orientation probabilities and leading evidence constants]
In the hard-margin posterior, assign each of its $K$ exceptional
neurons to the nearest critical ray and let $S_N$ be its orientation
set. Then
\[
 \Pr(S_N=S\mid H_N)\longrightarrow
 \frac{w_H(S)}{\sum_{T\in\mathcal B}w_H(T)},\qquad
 w_H(S)=\det(A_S^\top\Sigma^{-1}A_S)^{-1/2}.
 \tag{1}
\]
Every balanced pattern has a strictly positive limiting probability.
More precisely,
\[
 P_0^N(H_N)\sim e^{-NI(s)}N^{-Kr/2}
 L_r^K\alpha^{-K\nu}(M\rho)^M
 (2\pi)^{-(M-K)/2}(\det\Sigma)^{-1/2}
 \sum_{S\in\mathcal B}w_H(S).\tag{2}
\]
For the above-critical square or logistic likelihood, put
$s=s_\lambda$, $\alpha=(s-m_c)/(K\rho)$, and
$\gamma=\lambda\ell''(s)/M>0$. Then the orientation weights are
\[
 w_L(S)=\det\left[A_S^\top(I+\gamma\Sigma)^{-1}A_S\right]^{-1/2},
 \tag{3}
\]
and the likelihood evidence satisfies
\begin{align*}
 \mathcal Z_{N,\lambda}\sim{}&
 e^{-N[I(s)+\lambda\ell(s)]}N^{-K(r-2)/2}
 L_r^K\alpha^{-K\nu}(2\pi/\gamma)^{K/2}\\
 &\qquad\cdot\det(I+\gamma\Sigma)^{-1/2}
                  \sum_{S\in\mathcal B}w_L(S).\tag{4}
\end{align*}
No simultaneous dimension--width limit or effective sampling theorem is
asserted in this paper. The separate independent-construction companion
uses these constants to prove fixed-dimension sampling limits. The covariance and determinant expressions are exact finite-
dimension formulas, not numerically certified evaluations.
\end{theorem}

\section{The local central limit estimate actually needed}
Let $V=\Phi(\theta)$ under $Q_*$, with mean $\mu=m_c\chi$.
The companion constructs a fixed convolution block with a nonzero
centered isotropic Gaussian component. For $n$ neurons, separate the
block Bernoulli choices and discard the event with fewer than half their
expected Gaussian choices. Its probability is $O(e^{-cn})$.
The remaining subprobability law of $Y_n=\sum_{j=1}^nV_j$ has density
$p_n$ and is a mixture of
\[
 T+N(0,J\sigma^2 I),\qquad cn\le J\sigma^2\le Cn,
\]
where $T$ is the sum of the nongenuine-Gaussian block contributions and
leftovers. Its unconditional centered second moment is $O(n)$:
$\E T=n\mu$ and $\E\|T-n\mu\|^2\le Cn$.

The centered and rescaled good subprobability characteristic function
is bounded in absolute value by $e^{-c\|t\|^2}$. It converges pointwise
to $e^{-t^\top\Sigma t/2}$: the full law has the ordinary finite-variance
multivariate central limit theorem, and the discarded mass is
exponentially small. The CLT follows, for example, by expanding each
centered characteristic function to second order and taking its $n$th
power. Dominated Fourier inversion therefore proves
\[
 \sup_{z\in\R^M}
 \left|n^{M/2}p_n(n\mu+\sqrt n\,z)-\varphi_\Sigma(z)\right|
                 \longrightarrow0.\tag{5}
\]
Only the good density is used in (5). No assertion that the entire
finite convolution has a density is necessary.

A second consequence is tightness after integrating along any fixed
rank-$K$ subspace $E\subset\R^M$. The Gaussian component's density
integrated over $E$ is bounded by $Cn^{-(M-K)/2}$. If the residual
translation has centered norm above $B\sqrt n$, its contribution is
at most $Cn^{-(M-K)/2}/B^2$, by Chebyshev. On the complementary
residual event, Gaussian tails bound the part of the $E$ integral whose
coordinate has norm above $CB\sqrt n$ by
$Cn^{-(M-K)/2}e^{-cB^2}$, provided the remaining offset is bounded by
a smaller multiple of $B\sqrt n$.
These estimates are uniform in the mixture's translations and variances
on the stated events. They will justify extending bounded amplitude
windows to all of $\R^K$.

\section{One prescribed assignment of exceptional indices}
Fix $S\in\mathcal B$ and assign its ordered rays to $K$ distinct
neuron indices. Use the companion's coordinates
\[
 c=AR,\quad v=A-R,\quad t=\sqrt c\,u.
\]
Uniformly for bounded $v,t$ as $c\to\infty$, the one-neuron density is
\[
 \frac{(2\pi)^{-(r+2)/2}}{2M_*}
 c^{-\nu}\exp[-v^2/2-h_z(t)/\rho]\,(1+o(1))\,dc\,dv\,dt.
 \tag{6}
\]
This follows directly from the exact Jacobian
$c^{-r/2}/(A+R)$ and
$D_0=v^2/2+\|u\|^2/2+(A/\rho)h_z(u)$.
The output is $cf_z+e$ with $\|e\|_\infty\le C\|t\|$.
The uniform upper envelope is
$Cc^{-\nu}e^{-c_0(v^2+\|t\|)}$ in the fixed ray cone.

Put $n=N-K$ and
\[
 c_N=\frac{Ns-nm_c}{K\rho}=N\alpha+m_c/\rho,
 \qquad c_j=c_N+\sqrt N\,q_j.
\]
After summing the exceptional outputs, the remaining ordinary sum must
be close to $n\mu-\sqrt N A_Sq$ for hard margins. For bounded $q,v,t$
and bounded final slack, (5) applies uniformly; the bounded transverse
output error vanishes on the $\sqrt N$ scale.
Integrating (6) in $v,t$ gives the factor $L_r$ per ray. Its amplitude
scaling gives $N^{-r/2}\alpha^{-\nu}$ per ray.

For hard margins, the integral of the residual weight over the final
slack coordinates is exactly
\[
 \int_{(0,\infty)^M}
       \exp[-(M\rho)^{-1}\textstyle\sum_x y_x],dy=(M\rho)^M.
\]
Thus one assignment contributes asymptotically
\[
 N^{-Kr/2-M/2}L_r^K\alpha^{-K\nu}(M\rho)^M
                  \int_{\R^K}\varphi_\Sigma(A_Sq)\,dq.
 \tag{7}
\]
The Gaussian integral is
\[
 (2\pi)^{-(M-K)/2}(\det\Sigma)^{-1/2}
                     \det(A_S^\top\Sigma^{-1}A_S)^{-1/2}.
 \tag{8}
\]
For the soft likelihood, its centered residual weight at fluctuation
$z\sqrt N$ tends to $e^{-\gamma\|z\|^2/2}$. The analogous contribution
of one assignment is
\[
 N^{-Kr/2}L_r^K\alpha^{-K\nu}
       \int_{\R^K}\E_{Z\sim N(0,\Sigma)}
                  e^{-\gamma\|A_Sq+Z\|^2/2}\,dq.
 \tag{9}
\]
First integrating $Z$, then $q$, evaluates this integral as
\[
 (2\pi/\gamma)^{K/2}\det(I+\gamma\Sigma)^{-1/2}
       \det[A_S^\top(I+\gamma\Sigma)^{-1}A_S]^{-1/2}.\tag{10}
\]

Here are the domination details for the unbounded integrations in
(7)--(9). Restrict initially all $c_j$ to
$[N\alpha/2,2N\alpha]$, so $c_j^{-\nu}\le CN^{-\nu}$.
The exact-count volume bound, with the error-enlarged residual weight,
controls the $v,t$ tails by an integrable polynomial times
$\exp[-c_0\sum(v_j^2+\|t_j\|)]$ at the same power of $N$ as (7) or
(9). After fixing a large transverse box, use the subspace tightness
following (5) to control $q$ outside a large bounded set. For hard
margins first discard large final slack using its exponential tail.
For the soft likelihood put final slack equal to $\sqrt N z$;
its residual weights have an integrable bound
$C\exp[-c\min(\|z\|^2,\|z\|)]$ uniformly in $N$.
This follows from the quadratic minimum and positive linear growth of
$t_*y+\lambda\ell(y)$; square loss has the stronger global quadratic
bound. Split this $z$ integral into a bounded set and its uniformly
small tail, and apply the same subspace tightness on the bounded set.
Consequently the amplitude tails are $O(B^{-2})+O(e^{-cB})$ after
normalization, uniformly in $N$. The bad Gaussian-component mass is
exponentially small and hence negligible even after every polynomial
volume factor. These observations justify taking all bounding boxes
to infinity after the $N$ limit.

\section{Counting without hidden overcounting}
For each fixed $S$, there are $(N)_K=N(N-1)\cdots(N-K+1)$ assignments
of its ordered antipodal pairs to distinct indices. There is no extra
$K!$ divisor: the pair order is fixed, and the same network with exactly
one exceptional neuron in each pair has exactly one such assignment.

The exact-count theorem says that, with posterior probability tending
to one, precisely these $K$ neurons occur, their amplitudes divided by
$N$ approach $\alpha$, their angles approach their rays, and their
orientation set lies in $\mathcal B$. Hence the above assignment
events cover posterior mass $1-o(1)$. It remains to check that rare
overlap cannot change the constant. Above a fixed norm cutoff,
fourth-moment Markov and independence give
$Q_*^N(L_N>D)\le C_DN^{-D-1}$. Choose a fixed integer $D$ so large
that $N^K$ times this bound is negligible relative to both residual
normalizers from the exact-count theorem. On $L_N\le D$, assignment
events overlap at most a fixed number $(D)_K$, independently of $N$.
Their contribution with $L_N\ne K$ is therefore $o(1)$ of the residual
normalizer by the exact-count theorem. The contribution with $L_N>D$
is negligible by the preceding bound and $W_N\le1$.
Thus summing the one-assignment asymptotics gives the actual normalizer,
with relative error $o(1)$. Multiplication by $(N)_K\sim N^K$, using
$M=2K$, proves (2) and (4); the ratio for a fixed $S$ proves (1) and (3).

\section{Only a low-dimensional determinant distinguishes orientations}
The law of $\Phi$ under $Q_*$ is invariant under signed cube translations:
a coordinate sign change of the input is accompanied by the parity sign
change of the readout. Its covariance is therefore diagonal in the
orthonormal Walsh basis, and coordinate permutation symmetry makes its
eigenvalues depend only on Fourier degree. Let $D_E$ be its restriction
to the even Fourier subspace and let $\lambda_1>0$ be its common
singleton eigenvalue.

In that orthonormal basis write the columns of $A_S$ as an even block
$E$, a singleton block $L_S$, and zeros on all other odd degrees.
The even block is invertible and independent of the orientation choices
once a canonical representative of each antipodal pair is fixed. Put
$B=E^\top D_E^{-1}E$. The matrix determinant lemma gives
\[
 w_H(S)\ \propto\
 \det\left(I_r+\lambda_1^{-1}L_SB^{-1}L_S^\top\right)^{-1/2}.
\]
Thus the pattern-dependent part of a $K$-dimensional determinant reduces
to an $r$-dimensional one. For the soft case replace $D_E^{-1}$ by
$(I+\gamma D_E)^{-1}$ and $\lambda_1^{-1}$ by
$(1+\gamma\lambda_1)^{-1}$. This does not imply that all balanced
patterns have equal probabilities.

\section{Balanced orientations are resilient Boolean functions}
Choose canonical pair representatives $z=(1,b)$ with
$b\in\{-1,1\}^{r-1}$, and describe a selected orientation by
$\varepsilon(b)(1,b)$, where $\varepsilon$ is Boolean-valued.
Since $r$ is even, the balancing equations are exactly
\[
 \E_b g(b)=0,\qquad \E_b b_i g(b)=0\quad(1\le i\le r-1),
 \qquad g(b)=\chi_{r-1}(b)\varepsilon(b).
\]
Thus $g$ is a balanced first-order correlation-immune Boolean function,
also called a 1-resilient function. Equivalently, the real Fourier degree
of $\varepsilon$ is at most $r-3$; this is a statement about real Walsh
coefficients, not its algebraic degree over $\mathbb F_2$.
The combinatorial objects and their enumeration are established in the
cryptographic literature. For example,
\href{https://research.cs.wisc.edu/techreports/2007/TR1616.pdf}{the 2007
primary manuscript on counting correlation-immune Boolean functions}
records 807980 first-order resilient functions on five variables.
Our independent exact enumeration reproduces that count for $r=6$.
It splits the five-variable cube into two halves, enumerates their
$\binom{16}{8}$ balanced sign vectors, and pairs opposite four-coordinate
moment vectors. No floating-point test is needed for those conditions.

There is also an elementary lower bound at all dimensions. Every function
$g(b)=b_1h(b_2,\ldots,b_{r-1})$, with $h$ any balanced Boolean function,
is 1-resilient. Therefore
\[
 |\mathcal B|\ge\binom{K/2}{K/4}
                 \ge\frac{2^{K/2}}{K/2+1}.\tag{11}
\]
This weaker bound suffices to show that a fixed simple orientation family
has extremely small posterior probability.

\begin{theorem}[Simple hemisphere starts have negligible limiting mass]
In the hard-margin orientation law (1), uniformly over orientation
families $\mathcal A_r\subseteq\mathcal B$,
\[
 \Pr(S\in\mathcal A_r)\le |\mathcal A_r|\,
                        2^{-K/2+O(r^2)}.
\]
In particular the $2r$ coordinate hemispheres have probability at most
$2^{-2^{r-2}+O(r^2)}$ as even $r\to\infty$, after the width limit.
\end{theorem}
\begin{proof}
We bound the distortion introduced by the determinant weights. For
$r\ge8$ the critical law has $\E\|\theta\|^4=O(r^2)$.
Here are explicit sufficient estimates. Write $c=\sqrt{2/\pi}$ and
$T=\min_i|w_i|$. The minimum-half-normal calculation and Mills bound
from the critical-law companion give
\[
 \E_{Q_*}(G/\rho)^4
 \le \frac{24r^3}{c^4\prod_{j=2}^6(r-j)}=O(r^{-2}).
\]
Conditional Gaussian readouts give
$\E a^4=3+6\E(G/\rho)^2+\E(G/\rho)^4=O(1)$.
The incoming density is homogeneous with
$D=\|U\|^2-G^2/\rho^2$ in its exponent. Scaling its integral gives
$\E D^2=(r+1)(r+3)$, and hence $\E\|U\|^4=O(r^2)$.
Consequently
\[
 \|\Sigma\|\le\E\|\Phi\|^2
 \le M(r+1)\sqrt{\E a^4\E\|U\|^4}=O(Mr^2).
\]
For the singleton eigenvalue, conditional readout variance one and the
incoming density ratio $e^{G^2/(2\rho^2)}/M_*\ge1/M_*$ imply
$\lambda_1\ge M\E_0\widehat h_U(\{i\})^2/M_*$.
Gaussian integration by parts gives
$\E_0[w_i\widehat h_U(\{i\})]=1/2$; Cauchy--Schwarz therefore
bounds that square moment below by $1/4$. The rejection envelope in
the companion gives $M_*\le r/(r-2)$, so $\lambda_1\ge M(r-2)/(4r)$.

Every nonzero even coefficient of the critical feature has magnitude
at least $1/(M\sqrt r)$: its Fourier numerator is a nonzero integer
sum of positive parts of integer cube scores. In the reduced determinant
formula let $d_{\max}=\max_T\lambda_T/\widehat h(T)^2$, over even $T$.
Then $d_{\max}=O(M^3r^3)$. The singleton rows on the antipodal quotient
are orthogonal, and the even Walsh matrix is orthogonal up to its
normalization. These two facts give
\[
 0\preceq\lambda_1^{-1}L_SB^{-1}L_S^\top
       \preceq\frac{d_{\max}}{4r\lambda_1}I_r
       \preceq C M^2r^2I_r.
\]
Thus, after canceling the common factor $(\det B)^{-1/2}$, every
orientation weight lies between
$(1+CM^2r^2)^{-r/2}=2^{-O(r^2)}$ and one. Combine this with (11) to
obtain the claimed probability bound.
\end{proof}
This is a statement about posterior mass, not a lower bound on a specified
Markov chain's mixing time. It does show why remaining in a canonical
hemisphere sector cannot validate sampling of the full orientation law.

\paragraph{Six-bit numerical evaluation.}
Using an independently sampled critical covariance, the determinant
formula assigns approximately $1.073\times10^{-5}$ total probability
to the 12 coordinate hemispheres, compared with $12/807980\approx
1.485\times10^{-5}$ under the uniform orientation law. The minimum and
maximum reduced weights are approximately $0.163$ and $0.253$; the
large difference in family mass is primarily combinatorial.
These determinant values use floating-point Monte Carlo covariance
estimates and are not interval certificates. The enumeration count is
exact. The covariance sampler uses independent importance proposals
with a heavier tail than the earlier rejection sampler; its proposal
and numerical checks are recorded with the reproduction scripts.

\section{The remaining random part of the learned kernel}
The even-even block of the limiting feature kernel was already identified
in the companion. The singleton-singleton exceptional block is also
independent of $S$:
\[
 \E_{x,y}x_i[\mathcal K-\mathcal K_*](x,y)y_j
       =\frac{s-m_c}{4r\rho}\,1_{i=j}.
\]
Indeed the singleton coefficient of $h_z$ is $z_i/(2\sqrt r)$,
and summing $z_i z_j$ over one representative of each antipodal pair
gives $K1_{i=j}$. No exceptional feature has another odd Fourier degree.
Only the even-singleton cross block depends on the random orientation
pattern. For even $T$ its entry is
\[
 \frac{s-m_c}{K\rho}\frac{\widehat h(T)}{2\sqrt r}
       \sum_{z\in S}\chi_T(z)z_i.
\]
The parity row is zero by the balancing equations, which explains why
the target remains a deterministic kernel eigenvector. Differences
between two orientation-conditioned limiting kernels have rank at most
$2r$. Thus exponentially many exceptional directions coexist with only
a low-rank residual ambiguity in their feature kernel.

\section{Typical orientations have a small cross block}
A sharper existing enumeration theorem upgrades the preceding count into
a concentration result. This is useful because individual balanced
patterns can have very bad raw-feature transfer, as the finite-energy
transfer companion shows.

Let $n=r-1$ and $K=2^n$. Corollary 1.1 of Canfield, Gao, Greenhill,
McKay, and Robinson, \emph{Asymptotic enumeration of correlation-immune
boolean functions} (2010), gives, at order one,
\[
 |\mathcal B|\sim 2^{K+n}(\pi 2^{n-1})^{-r/2}.
 \tag{12}
\]
See the \href{https://web.maths.unsw.edu.au/~csg/papers/immune.pdf}{authors'
primary manuscript}. In particular, independent uniform orientation
signs satisfy the balancing equations with probability at least
$e^{-C_1r^2}$ for all sufficiently large even $r$. This counting theorem
is an external input; the exact enumeration of five-variable functions
is not its proof.

For an odd set $B\subseteq[r]$ define
\[
 m_B(S)=\frac1K\sum_{z\in S}\chi_B(z).
\]
In the canonical description $z=\varepsilon(b)(1,b)$, this is a Walsh
coefficient of the independent signs $\varepsilon(b)$ before balancing.
The $K$ odd characters give exactly the $K$ distinct characters of $b$.
Hoeffding and a union bound therefore give
\[
 \Pr_{\rm independent}\{\max_{|B|\text{ odd}}|m_B|>t\}
                              \le2K e^{-Kt^2/2}.
\]
The determinant distortion bound already proved above is $e^{O(r^2)}$.
Dividing by the balancing probability from (12) thus proves the following.

\begin{theorem}[Orientation Fourier concentration]
There are absolute constants $C,c>0$ such that the hard-margin limiting
orientation law satisfies, for all sufficiently large even $r$,
\[
 \Pr\left\{\max_{|B|\text{ odd}}|m_B(S)|>
                           \frac{Cr}{\sqrt K}\right\}
                                                    \le e^{-cr^2}.
 \tag{13}
\]
More explicitly its tail is bounded by
$2K\exp(C_2r^2-Kt^2/2)$ for a fixed constant $C_2$.
\end{theorem}
\begin{proof}
For any event of orientation signs, its posterior probability is at most
its independent-sign probability multiplied by the reciprocal balancing
probability and the ratio of the largest and smallest determinant
weights. The preceding two bounds give $\exp(C_2r^2)$ for that multiplier.
Take $t=Cr/\sqrt K$ with $C$ sufficiently large. The factor $2K$ has
logarithm $O(r)$ and is absorbed in the negative quadratic exponent.
\end{proof}

This does not assert that the orientation law is close to uniform in
total variation. It bounds its Walsh coefficients despite its
nonuniform determinant weights. The raw-feature transfer companion uses
(13) to show that a typical exceptional-neuron representation has only
an exponentially small linear error on every nearby parity task.

\section{Scope and open work}
This note provides formulas and an analytic asymptotic argument. It does
not prove that a practical chain mixes over the orientation sets, supply
interval-certified covariance values, or classify the most probable sets
as $r$ grows. The six-bit evaluation is exploratory. The $r=6$ critical law has finite output
covariance but need not have a finite variance for naive covariance
estimators; ordinary iid error bars for those entries would require
additional care. The earlier independent sampler estimates $m_c$ with
finite variance, which is a different observable.
\end{document}
